\section{Math (Ruiming)}

\entry{exgcd.cpp}{Ruiming, verbatim}{Extended Euclid; prints the inverse of \code{a} mod \code{b}.}

%NOCOMPILE
%VERBATIM
\begin{lstlisting}[breaklines=true]
// header lines: see NewCodeforces.cpp, geomtery/Line_and_Vectors.cpp
ll n,a[MAXN],ans;
ll x,y,px,py,d;
void exgcd(ll a,ll b){
    if(b==0){
        x=1;
        y=0;
        d=a;    
    }else{
        exgcd(b,a-a/b*b);
        px=x;
        py=y;
        x=py;
        y=px-py*(a/b);
    }
}
int main(){
    ios_base::sync_with_stdio(false);
    cin.tie(NULL);
    ll a,b;
    cin>>a>>b;
    exgcd(a,b);
    cout<<(x%b+b)%b<<'\n';
}
\end{lstlisting}

\entry{CRT.cpp}{Ruiming, verbatim}{Chinese remainder theorem by stepping through residues (small coprime moduli).}

%NOCOMPILE
%VERBATIM
\begin{lstlisting}[breaklines=true]
// header lines: see NewCodeforces.cpp, geomtery/Line_and_Vectors.cpp
const ll MAXN=16;//UPDATERA ARRAY STORLEKEN!!!!!!!!!!!!!!!!!!!!!!!!!!!!!!!!!!!!
ll n,lst[MAXN],a[MAXN],m[MAXN],ans[MAXN],k,temp;
int main(){
    ios_base::sync_with_stdio(false);
    cin.tie(NULL);
    cin>>n;
    for(int i=1;i<=n;i++){
        cin>>m[i]>>a[i];
    }
    ans[1]=a[1];
    for(int i=2;i<=n;i++){
        for(ll j=0;j<=m[i]-1;j++){
            if(((j*m[i-1]+ans[i-1])%m[i])==a[i]){
                ans[i]=j*m[i-1]+ans[i-1];
                break;
            }            
        }
        m[i]*=m[i-1];
    }
    cout<<ans[n]<<'\n';
}
\end{lstlisting}

\entry{Fast\_gcd.cpp}{Ruiming, verbatim}{Binary gcd with \code{\_\_builtin\_ctz}.}

%NOCOMPILE
%VERBATIM
\begin{lstlisting}[breaklines=true]
// header lines: see NewCodeforces.cpp
ll gcd(ll a,ll b)
{
	ll az=__builtin_ctz(a),bz=__builtin_ctz(b);
	ll z=min(az,bz);
	ll dif;
	b>>=bz;
	while(a)
	{
		a>>=az;
		dif=b-a;
		az=__builtin_ctz(dif);
		if(a<b) b=a;
		if(dif<0) a=-dif;
		else a=dif;
	}
	return b<<z;
}
\end{lstlisting}

\entry{Fastmod.cpp}{Ruiming, verbatim}{\code{mul} (safe \code{a*b\%m}) and \code{fast} ($x^p \bmod m$).}

%NOCOMPILE
%VERBATIM
\begin{lstlisting}[breaklines=true]
// header lines: see NewCodeforces.cpp, geomtery/Line_and_Vectors.cpp, math/exgcd.cpp
ll mul(ll a, ll b, ll m) { return (__int128)a * b % m; }
ll fast(ll x,ll p,ll m){
    x=(x%m+m)%m;
    ll ans=1%m;
    while(p>0){
        if(p&1){
            ans=mul(ans,x,m);
        }
        x=mul(x,x,m);
        p>>=1;
    }
    return ans;
}
int main(){
    ios_base::sync_with_stdio(false);
    cin.tie(NULL);
    cin>>n;
}
\end{lstlisting}

\entry{Eulerthm.cpp}{Ruiming, verbatim}{Euler's theorem: $a^b \bmod m$ for a huge exponent $b$.}

%NOCOMPILE
%VERBATIM
\begin{lstlisting}[breaklines=true]
// header lines: see NewCodeforces.cpp
ll mul(ll a, ll b, ll m) { /* see math/Fastmod.cpp */ }
ll fast(ll x,ll p,ll m){ /* see math/Fastmod.cpp */ }
bool flag;
ll gcd(ll x,ll y){
    while(x!=y){
        if(x>y) swap(x,y);
        if(y%x==0) return x;
        y-=x*(y/x);
    }
    return x;
}
ll a,m,phi=1;
string b;
ll read(){
    string x;
    cin>>x;
    ll ans=0;
    for(int i=0;i<x.size();i++){
        ans=(ans*10+(x[i]-'0'));
        if(ans>=phi) flag=true;
        ans%=phi;
    }
    return ans;
}
int main(){
    cin>>a>>m;
    //find the euler function to m
    ll d=gcd(a,m);
    ll temp=m;
    for(int i=2;i<=temp;i++){
        if(temp%i==0){//here uses temp still makes sure i is a prime
            ll times=0;
            while(temp%i==0){
                times++;
                temp/=i;//makes sure i is a prime
            }
            phi*=pow(i,times-1)*(i-1);
        }
    }
    //if(phi==1) phi=m-1;
    //b needs to be stored as a string and takes mod phi
    //then we can calculate it with fast mod
    ll b=read();
    ll ans;
    if(flag) ans=fast(a,b+phi,m);
    else ans=fast(a,b,m);
    //deal with d d|m, will lead up to a , should be some period
    cout<<ans<<'\n';
}
\end{lstlisting}

\entry{Phi.cpp}{Ruiming, verbatim}{Euler's $\varphi(n)$ for one $n$.}

%NOCOMPILE
%VERBATIM
\begin{lstlisting}[breaklines=true]
// header lines: see NewCodeforces.cpp, geomtery/Line_and_Vectors.cpp, math/exgcd.cpp
ll phi(ll n) {//求欧拉函数
    ll ans = 1;
    for(ll i = 2; i * i <= n; ++ i) {
        if(n % i == 0) {
            n /= i;
            ans *= i - 1;
            while(n % i == 0) {
                n /= i;
                ans *= i;
            }
        }
    }
    if(n > 1)//优化，n>1表示分解剩下一个质数
        ans *= n - 1;
    return ans;
}
int main(){ /* see math/Fastmod.cpp */ }
\end{lstlisting}

\entry{Linearphi.cpp}{Ruiming, verbatim}{$\varphi$ for all numbers up to $n$ (linear sieve).}

%NOCOMPILE
%VERBATIM
\begin{lstlisting}[breaklines=true]
// header lines: see NewCodeforces.cpp, geomtery/Line_and_Vectors.cpp
const ll MAXN=406;//UPDATERA ARRAY STORLEKEN!!!!!!!!!!!!!!!!!!!!!!!!!!!!!!!!!!!!
ll n,phi[MAXN],ans=1;
bool p[MAXN];
vector<ll> primes;
int main(){
    ios_base::sync_with_stdio(false);
    cin.tie(NULL);
    cin>>n;
    if(n==1) cout<<0<<'\n';
    else{
        phi[1]=1;
        for(int i=2;i<=n;i++){//modiying euler sieve
            if(!p[i]){
                primes.push_back(i);
                phi[i]=i-1;
            }
            for(int j=0;j<primes.size();j++){
                if(primes[j]*i>n) break;
                p[primes[j]*i]=1;
                if(i%primes[j]==0){
                    phi[i*primes[j]]=primes[j]*phi[i];
                    break;
                }else{
                    phi[i*primes[j]]=(primes[j]-1)*phi[i];
                }
            }
        }
    }
}
\end{lstlisting}

\entry{Euler\_sieve.cpp}{Ruiming, verbatim}{Linear (Euler) sieve of primes.}

%NOCOMPILE
%VERBATIM
\begin{lstlisting}[breaklines=true]
// header lines: see NewCodeforces.cpp, math/Linearphi.cpp
const int MAXN=106;
ll n,q;
int main(){
    ios_base::sync_with_stdio(false);
    cin.tie(NULL);
    cin>>n>>q;
    for(ll i=2;i<=n;i++){
        if(!p[i]){
            primes.push_back(i);
        }
        for(ll j=0;j<primes.size();j++){
            if(primes[j]*i>n) break;
            p[primes[j]*i]=1;
            if(i%primes[j]==0) break;
        }
    }
    for(int i=1;i<=q;i++){
        ll a;
        cin>>a;
        cout<<primes[a-1]<<'\n';
    }
}
\end{lstlisting}

\entry{Factorization.cpp}{Ruiming, verbatim}{Factorises every number below \code{MAXN} by trial division into \code{map}s.}

%NOCOMPILE
%VERBATIM
\begin{lstlisting}[breaklines=true]
// header lines: see NewCodeforces.cpp, geomtery/Line_and_Vectors.cpp, math/exgcd.cpp
const ll MAXN=36;//UPDATERA ARRAY STORLEKEN!!!!!!!!!!!!!!!!!!!!!!!!!!!!!!!!!!!!
map<ll,ll> fac[MAXN];
int main(){
    ios_base::sync_with_stdio(false);
    cin.tie(NULL);
    cin>>n;
    for(int i=1;i<MAXN;i++){
        ll temp=i;
        for(int j=2;j<=sqrt(temp+0.5);j++){
            while(temp%j==0){
                fac[i][j]++;
                temp/=j;
            }
        }
        if(temp>1){
            fac[i][temp]++;
        }
    }
    for(int i=1;i<=n;i++){
        cin>>a[i];
    }
}
\end{lstlisting}

\entry{Prime\_factorization.cpp}{Ruiming, verbatim}{Prime factorisation of one number, \code{factor(n)}.}

%NOCOMPILE
%VERBATIM
\begin{lstlisting}[breaklines=true]
// header lines: see NewCodeforces.cpp
vector<ll> factor(ll n) {
	vector<ll> ret;
	for (ll i = 2; i * i <= n; i++) {
		while (n % i == 0) {
			ret.push_back(i);
			n /= i;
		}
	}
	if (n > 1) { ret.push_back(n); }
	return ret;
}
const ll INF=1e18;
const ll MAXN=206;//UPDATERA ARRAY STORLEKEN!!!!!!!!!!!!!!!!!!!!!!!!!!!!!!!!!!!!
ll n,a[MAXN],ans;
int main(){
    ios_base::sync_with_stdio(false);
    cin.tie(NULL);
    cin>>n;
    vector<ll> fac;
    fac=factor(n);
}
\end{lstlisting}

\entry{miller\_rabin.cpp}{Ruiming, verbatim}{Miller-Rabin primality test.}

%NOCOMPILE
%VERBATIM
\begin{lstlisting}[breaklines=true]
// header lines: see NewCodeforces.cpp
typedef __int128 ll;
typedef pair<int, int> pii;
template<typename T> inline T read() {
    T x = 0, f = 1; char ch = 0;
    for(; !isdigit(ch); ch = getchar()) if(ch == '-') f = -1;
    for(; isdigit(ch); ch = getchar()) x = (x << 3) + (x << 1) + (ch - '0');
    return x * f;
}
template<typename T> inline void write(T x) {
    if(x < 0) putchar('-'), x = -x;
    if(x > 9) write(x / 10);
    putchar(x % 10 + '0');
}
template<typename T> inline void print(T x, char ed = '\n') {
    write(x), putchar(ed);
}
ll t, n;
ll mul(ll a, ll b, ll m) { /* see math/Fastmod.cpp */ }
ll fast(ll x,ll p,ll m){ /* see math/Fastmod.cpp */ }
vector<ll> p = {2, 3, 5, 7, 11, 13, 17, 19, 23, 29, 31, 37};
bool miller_rabin(ll n) {
    if(n < 3 || n % 2 == 0) return n == 2;
    ll u = n - 1, t = 0;
    while(u % 2 == 0) u /= 2, ++ t;
    for(auto a : p) {
        if(n == a) return 1;
        if(n % a == 0) return 0;
        ll v = fast(a, u, n);
        if(v == 1) continue;
        ll s = 1;
        for(; s <= t; ++ s) {
            if(v == n - 1) break;
            v = v * v % n;
        }
        if(s > t) return 0; 
    }
    return 1;
}
int main() {
    t = read<ll>();
    while(t --) {
        n = read<ll>();
        if(miller_rabin(n)) puts("Yes");
        else puts("No");
    }
    return 0;
}
\end{lstlisting}

\entry{Pollard\_Rho.cpp}{Ruiming, verbatim}{Pollard rho factorisation for numbers up to $\sim 10^{18}$.}

%NOCOMPILE
%VERBATIM
\begin{lstlisting}[breaklines=true]
// header lines: see geomtery/Half_Plane_Intersect.cpp, NewCodeforces.cpp
#define LL long long 
#define i128 __int128
mt19937_64 rd(time(0));
LL n;
LL mul(LL a, LL b, LL m) { return (__int128)a * b % m; }
LL fast(LL x,LL p,LL m){
    x=(x%m+m)%m;
    LL ans=1%m;
    while(p>0){
        if(p&1){
            ans=mul(ans,x,m);
        }
        x=mul(x,x,m);
        p>>=1;
    }
    return ans;
}
bool miller_rabin(LL n)
{
    if(n <= 1 || n > 2 && ~n & 1) return false;
    LL m = n - 1, k = 0;
    while(~m & 1) m >>= 1, k++;
    vector<int> s = {2, 3, 5, 7, 11, 13, 17, 19, 23, 29, 31, 37};
    for(int a : s)
    {
        if(a >= n) break;
        i128 x = fast(a, m, n), y;
        for(int i = 1; i <= k; i++, x = y)
        {
            y = x * x % n;
            if(y == 1 && x != 1 && x != n - 1) return false;
        }
        if(x != 1) return false;
    }
    return true;
}
vector<LL> p;
void pollard_rho(LL n)
{
	if(n <= 1) return;
    if(miller_rabin(n))
    {
        p.push_back(n);
        return;
    }
    LL d = rd() % (n - 1) + 1;
    auto f = [&](LL x)
    {
        return ((i128)x * x + d) % n; 
    };
    LL x = rd() % (n - 1) + 1, y = f(x), t = 1;
    while(t == 1)
    {
        for(int i = 1; i <= 32; i++)
        {
            while(x == y)
                d = rd() % (n - 1) + 1, x = rd() % (n - 1) + 1, y = f(x);
            if((i128)t * abs(x - y) % n == 0) break;
            t = (i128)t * abs(x - y) % n;
            x = f(x), y = f(f(y));
        }
        t = __gcd(n, t);
    }
    pollard_rho(t), pollard_rho(n / t);
}
int main()
{
    int T;
    cin >> T;
    while(T--)
    {
        cin >> n;
        if(miller_rabin(n))
        {
            puts("Prime");
            continue;
        }
        p.clear();
        pollard_rho(n);
        LL ans = 0;
        for(LL i : p) ans = max(ans, i);
        cout << ans << endl;
    }
    return 0;
}
\end{lstlisting}

\entry{Combinumber.cpp}{Ruiming, verbatim}{Pascal's triangle \code{c[i][j]} mod \code{M}.}

%NOCOMPILE
%VERBATIM
\begin{lstlisting}[breaklines=true]
// header lines: see NewCodeforces.cpp, geomtery/Line_and_Vectors.cpp
typedef vector<ll> vi;
typedef vector<vi> matrix;
#define rep(i, a, b) for(int i = a; i < b; i++)
const ll M=1000000007;
ll c[MAXN][106],t,n,m,k,lst[MAXN],ans;
int main(){
    ios_base::sync_with_stdio(false);
    cin.tie(NULL);
    c[0][0]=1;
    for(int i=1;i<MAXN-1;i++){
        c[i][0]=1;
        if(i<106) c[i][i]=1;
        for(int j=1;j<min(i,105);j++){
            c[i][j]=c[i-1][j]+c[i-1][j-1];
            c[i][j]%=M;
        }
    }
}
\end{lstlisting}

\entry{Combinumber\_inv.cpp}{Ruiming, verbatim}{Binomial coefficients with modular inverses.}

%NOCOMPILE
%VERBATIM
\begin{lstlisting}[breaklines=true]
// header lines: see NewCodeforces.cpp, geomtery/Line_and_Vectors.cpp
const ll MOD=1e9+7;
const ll MAXN=26;//UPDATERA ARRAY STORLEKEN!!!!!!!!!!!!!!!!!!!!!!!!!!!!!!!!!!!!
ll t,n,fac[MAXN],inv[MAXN],ans,k;
ll mul(ll a, ll b, ll m) { /* see math/Fastmod.cpp */ }
ll fast(ll x,ll p,ll m){ /* see math/Fastmod.cpp */ }
int main(){
    ios_base::sync_with_stdio(false);
    cin.tie(NULL);
    cin>>n>>k;
    ll temp=1;
    for(int i=1;i<=n;i++){
        temp*=i;
        temp%=MOD;
        inv[i]=fast(temp,MOD-2,MOD);
        fac[i]=temp;
    }
    fac[0]=1;
    inv[0]=1;
    for(int i=0;i<=n;i++){
        ll comb=((fac[n]*inv[i])%MOD*inv[n-i])%MOD;
    }
}
\end{lstlisting}

\entry{Linearinverse.cpp}{Ruiming, verbatim}{Inverses of $1..n$ modulo a prime in $O(n)$.}

%NOCOMPILE
%VERBATIM
\begin{lstlisting}[breaklines=true]
// header lines: see NewCodeforces.cpp
const int MAXN=206;
ll n,p,ans[MAXN],corr[MAXN],now,fac[MAXN],inv[MAXN];
ll mul(ll a, ll b, ll m) { /* see math/Fastmod.cpp */ }
ll fast(ll x,ll p,ll m){ /* see math/Fastmod.cpp */ }
int main(){
    cin>>n>>p;
    fac[0]=1;
    for(int i=1;i<=n;i++){
        fac[i]=(fac[i-1]*i)%p;
    }
    inv[n]=fast(fac[n],p-2,p);//inv fac[i]
    for(int i=n-1;i>=1;i--){
        inv[i]=((i+1)*inv[i+1])%p;
    }
    for(int i=1;i<=n;i++){
        cout<<(fac[i-1]*inv[i])%p<<'\n';
    }
}
\end{lstlisting}

\entry{Matrix.cpp}{Ruiming, verbatim}{Matrix multiplication and fast matrix power mod \code{MOD}.}

%NOCOMPILE
%VERBATIM
\begin{lstlisting}[breaklines=true]
// header lines: see NewCodeforces.cpp, math/Combinumber_inv.cpp, geomtery/Line_and_Vectors.cpp
const ll MAXN=106;//UPDATERA ARRAY STORLEKEN!!!!!!!!!!!!!!!!!!!!!!!!!!!!!!!!!!!!
ll n,ans,k;
struct matrix{
    ll row,col;
    ll m[MAXN][MAXN];
};
matrix operator * (const matrix &a,const matrix &b){
    //assuming they fits a.col==b.row
    matrix c;
    memset(c.m,0,sizeof(c.m));
    c.row=a.row;
    c.col=b.col;
    for(int i=1;i<=c.row;i++){
        for(int j=1;j<=c.col;j++){
            for(int k=1;k<=a.col;k++){
                c.m[i][j]+=(a.m[i][k]*b.m[k][j])%MOD;
                c.m[i][j]%=MOD;
            }
        }
    }
    return c;
}
matrix pow_matrix(matrix a,ll n){
    //assuming a is a square matrix
    matrix res;
    memset(res.m,0,sizeof(res.m));
    res.row=res.col=a.col;
    for(int i=1;i<MAXN;i++) res.m[i][i]=1;
    while(n){
        if(n&1) res=res*a;
        a=a*a;
        n/=2;
    }
    return res;
}
int main(){
    ios_base::sync_with_stdio(false);
    cin.tie(NULL);
    cin>>n>>k;
    matrix A;
    A.row=A.col=n;
    for(int i=1;i<=n;i++){
        for(int j=1;j<=n;j++){
            cin>>A.m[i][j];
        }
    }
    matrix res=pow_matrix(A,k);
    for(int i=1;i<=n;i++){
        for(int j=1;j<=n;j++){
            cout<<res.m[i][j]<<' ';
        }
        cout<<'\n';
    }
}
\end{lstlisting}

\entry{gaussjordanelimination.cpp}{Ruiming, verbatim}{Gauss-Jordan elimination over the reals.}

%NOCOMPILE
%VERBATIM
\begin{lstlisting}[breaklines=true]
// header lines: see NewCodeforces.cpp, math/Combinumber_inv.cpp, geomtery/Line_and_Vectors.cpp, math/Matrix.cpp
ld eps=1e-7;//MIGHT GET WA PGA EPS TOO SMALL!!!!!!!!!!!!!!!!!!!!!!!!!!!!
ll ldcmp(ld x,ld y){ /* see geomtery/Line_and_Vectors.cpp */ }
struct matrix{
    ld row,col;
    ld m[MAXN][MAXN];
};
ld res[MAXN];
//A merged matrix
void gauss_elimination(matrix& A,ld *res){
    memset(res,0,sizeof(res));
    ll n=A.row;
    if(A.col!=n+1){
        for(int i=1;i<=n;i++){
            res[i]=INF;
        }
        return;
    }
    for(int i=1;i<=n;i++){
        ll maxi=i;
        //find pivot row
        for(int j=1;j<=n;j++){
            if(j<i&&ldcmp(A.m[j][j],0)>0){
                continue;
            }
            if(ldcmp(abs(A.m[j][i]),abs(A.m[maxi][i]))==1){
                maxi=j;
            }
        }
        swap(A.m[maxi],A.m[i]);
        if(ldcmp(A.m[i][i],0)!=0){
            ld tmp=A.m[i][i];
            for(int j=i;j<=n+1;j++){
                A.m[i][j]/=tmp;
            }
            for(int j=1;j<=n;j++){
                if(i!=j){
                    ld rate=A.m[j][i]/A.m[i][i];
                    for(int k=i;k<=n+1;k++){
                        A.m[j][k]-=A.m[i][k]*rate;
                    }
                }
            }
        }
    }
    return;
}
int main(){
    ios_base::sync_with_stdio(false);
    cin.tie(NULL);
    cin>>n;
    matrix A;
    A.row=n;
    A.col=n+1;
    for(int i=1;i<=n;i++){
        for(int j=1;j<=n+1;j++){
            cin>>A.m[i][j];
        }
    }
    gauss_elimination(A,res);
    ll sign=1;
    for(int i=1;i<=n;i++){
        if(ldcmp(abs(A.m[i][i]),0)==0&&ldcmp(abs(A.m[i][n+1]),0)!=0){
            sign=-1;
            break;
        }
        if(ldcmp(abs(A.m[i][i]),0)==0){
            sign=0;
        }
    }
    if(sign==1){
        cout<<fixed<<setprecision(2);
        for(int i=1;i<=n;i++){
            cout<<"x"<<i<<"="<<A.m[i][n+1]<<'\n';
        }
    }else cout<<sign<<'\n';
}
\end{lstlisting}

\entry{mod\_elimination.cpp}{Ruiming, verbatim}{Gaussian elimination modulo a prime.}

%NOCOMPILE
%VERBATIM
\begin{lstlisting}[breaklines=true]
// header lines: see NewCodeforces.cpp, geomtery/Line_and_Vectors.cpp
ll n,ans,k,MOD;
struct matrix{ /* see math/Matrix.cpp */ };
ll res[MAXN],inv[MAXN];
void lin_inv(){
    inv[1]=1;
    for(int i=2;i<MOD;i++){
        inv[i]=(MOD-inv[MOD%i]*(MOD/i)%MOD);
    }
}
//A merged matrix
void gauss_elimination(matrix& A,ll *res){
    memset(res,0,sizeof(res));
    ll n=A.row;
    if(A.col!=n+1){
        for(int i=1;i<=n;i++){
            res[i]=INF;
        }
        return;
    }
    for(int i=1;i<=n;i++){
        if(A.m[i][i]==0){
            for(int j=1;j<=n;j++){
                //even before not pivot row
                if(j<i&&A.m[j][j]!=0) continue;
                if(A.m[j][i]!=0){
                    swap(A.m[i],A.m[j]);
                    break;
                }
            }
        }
        if(A.m[i][i]!=0){
            for(int j=1;j<=n;j++){
                if(i!=j&&A.m[j][i]!=0){
                    ll d=__gcd(A.m[i][i],A.m[j][i]);
                    ll a=A.m[i][i]/d,b=A.m[j][i]/d;
                    for(int k=1;k<=n+1;k++){
                        A.m[i][k]*=b;
                        A.m[i][k]%=MOD;
                        A.m[j][k]*=a;
                        A.m[j][k]%=MOD;
                    }
                    for(int k=i;k<=n+1;k++){
                        A.m[j][k]-=A.m[i][k];
                        A.m[j][k]=(A.m[j][k]+MOD)%MOD;
                    }
                }
            }
        }
    }
    lin_inv();
    for(int i=1;i<=n;i++){
        if(A.m[i][i]!=0){
            bool flag=0;
            for(int j=i+1;j<=n;j++){
                if(A.m[i][j]!=0){
                    flag=1;
                    break;
                }
            }
            if(!flag){
                A.m[i][n+1]=(A.m[i][n+1]*inv[A.m[i][i]])%MOD;
                A.m[i][i]=1;
            }
        }
    }
    return;
}
int main(){
    ios_base::sync_with_stdio(false);
    cin.tie(NULL);
    cin>>n>>MOD;
    matrix A;
    A.row=n;
    A.col=n+1;
    for(int i=1;i<=n;i++){
        for(int j=1;j<=n+1;j++){
            cin>>A.m[i][j];
        }
    }
    gauss_elimination(A,res);
    for(int i=1;i<=n;i++){
        for(int j=1;j<=n+1;j++) cout<<A.m[i][j]<<' ';
        cout<<'\n';
    }
}
\end{lstlisting}

\entry{xor\_gauss.cpp}{Ruiming, verbatim}{Gaussian elimination over XOR (bitsets).}

%NOCOMPILE
%VERBATIM
\begin{lstlisting}[breaklines=true]
// header lines: see NewCodeforces.cpp, math/Combinumber_inv.cpp, geomtery/Line_and_Vectors.cpp
const ll MAXN=2006;//UPDATERA ARRAY STORLEKEN!!!!!!!!!!!!!!!!!!!!!!!!!!!!!!!!!!!!
ll n,ans,k,m;
struct matrix{
    ll row,col;
    bitset<MAXN> m[MAXN]={};
};
//A bit matrix 0/1
void xor_gauss(matrix& A){
    ll m=A.row;
    ll n=A.col;
    for(int i=1;i<=n-1;i++){
        ll maxi=i;
        //find pivot row
        for(int j=1;j<=m;j++){
            if(j<i&&A.m[j][j]>0){
                continue;
            }
            if(A.m[j][i]==1){
                swap(A.m[j],A.m[i]);
                k=max(k,(ll)j);
                break;
            }
        }
        if(A.m[i][i]==1){
            for(int j=1;j<=m;j++){
                if(i!=j&&A.m[j][i]==1){
                    A.m[j]^=A.m[i];
                }
            }
        }else{
            k=m+1;
            return;
        }
    }
    return;
}
int main(){
    ios_base::sync_with_stdio(false);
    cin.tie(NULL);
    cin>>n>>m;
    matrix A;
    A.row=m;
    A.col=n+1;
    for(int i=1;i<=m;i++){
        string s;
        cin>>s;
        for(int j=1;j<=n;j++){
            if(s[j-1]=='1') A.m[i][j]=1;
        }
        ll x;
        cin>>x;
        if(x) A.m[i][n+1]=1;
    }
    xor_gauss(A);
    if(k==m+1){
        cout<<"Cannot Determine"<<'\n';
    }else{
        cout<<k<<'\n';
        for(int i=1;i<=n;i++){
            if(A.m[i][n+1]==1){
                cout<<"?y7M#"<<'\n';
            }else cout<<"Earth"<<'\n';
        }
    }
}
\end{lstlisting}

\entry{xor\_basis.cpp}{Ruiming, verbatim}{XOR basis (linear basis); prints the maximum subset XOR.}

%NOCOMPILE
%VERBATIM
\begin{lstlisting}[breaklines=true]
// header lines: see NewCodeforces.cpp, geomtery/Line_and_Vectors.cpp
ll n,a[MAXN],ans,len=1,bits=64;
bool zero;
void get_basis(){
    for(int i=bits;i>=0;i--){
        for(int j=len;j<=n;j++){
            if(a[j]&(1LL<<(i-1))){
                swap(a[len],a[j]);
                break;
            }
        }
        if(a[len]&(1LL<<(i-1))){
            for(int j=1;j<=n;j++){
                if(j==len) continue;
                if(a[j]&(1LL<<(i-1))){
                    a[j]^=a[len];
                }
            }
            len++;
        }
    }
    len--;
    zero=(len!=n);
}
int main(){
    ios_base::sync_with_stdio(false);
    cin.tie(NULL);
    cin>>n;
    for(int i=1;i<=n;i++) cin>>a[i];
    get_basis();
    sort(a+1,a+len+1);
    for(int i=1;i<=len;i++) ans^=a[i];
    cout<<ans<<'\n';
}
\end{lstlisting}

\entry{zero\_one\_partition.cpp}{Ruiming, verbatim}{0/1 fractional programming (binary search on the ratio).}

%NOCOMPILE
%VERBATIM
\begin{lstlisting}[breaklines=true]
// header lines: see NewCodeforces.cpp, geomtery/Line_and_Vectors.cpp, geomtery/min_cover_circle.cpp
ll n,a[MAXN],b[MAXN],ans,k;
ld eps=1e-6;
bool check(ld x){
    vector<ld> all;
    ld res=0;
    for(int i=1;i<=n;i++){
        all.push_back(a[i]-x*b[i]);
    }
    sort(all.begin(),all.end(),greater<ld>());
    for(int i=0;i<n-k;i++){
        res+=all[i];
    }
    return res>=-eps;
}
ll zero_one_partition(){
    ld start=0,end=n;
    while(end-start>eps){
        ld x=start+(end-start)/2;
        //if f(x)>0
        if(check(x)){
            start=x+eps;
        }else end=x-eps;
    }
    ll res = floor(100 * start + 0.5);
    return res;
}
int main(){
    ios_base::sync_with_stdio(false);
    cin.tie(NULL);
    while(true){
        cin>>n>>k;
        if(n==0&&k==0) break;
        for(int i=1;i<=n;i++) cin>>a[i];
        for(int i=1;i<=n;i++) cin>>b[i];
        cout<<zero_one_partition()<<'\n';
    }
}
\end{lstlisting}

